\documentclass[../main]{subfiles}
\begin{document}

\chapter{Primer Principio de la Termodinámica}

\section{Primer Principio de la Termodinámica}

\subsection{Teoría}

El Primer Principio de la Termodinámica, también conocido como el principio de conservación de la energía para sistemas termodinámicos, establece que la energía no puede crearse ni destruirse, solo puede transferirse o transformarse de una forma a otra. Este principio es una extensión del principio de conservación de la energía a los procesos termodinámicos y es fundamental en el estudio de la termodinámica.

Matemáticamente, el Primer Principio de la Termodinámica se expresa como:
\begin{equation}
\Delta U = Q - W
\end{equation}
donde:
\begin{itemize}
    \item $\Delta U$: Cambio en la energía interna del sistema [J]
    \item $Q$: Calor añadido al sistema [J]
    \item $W$: Trabajo realizado por el sistema [J]
\end{itemize}

La energía interna ($U$) de un sistema es una función de estado que depende de la temperatura, el volumen y la cantidad de sustancia del sistema. El calor ($Q$) es la transferencia de energía debido a una diferencia de temperatura entre el sistema y su entorno. El trabajo ($W$) es la transferencia de energía que se realiza cuando una fuerza actúa sobre una distancia.

En términos de un ciclo termodinámico, el Primer Principio puede reformularse para establecer que la suma de las transferencias de calor y trabajo alrededor de un ciclo completo es cero:
\begin{equation}
\oint (dQ - dW) = 0
\end{equation}

Otra expresión del Primer Principio de la Termodinámica es a través de la capacidad calorífica a volumen constante ($C_V$) y la capacidad calorífica a presión constante ($C_P$). Para un gas ideal, estas se definen como:
\begin{equation}
C_V = \left( \frac{\partial U}{\partial T} \right)_V
\end{equation}
y
\begin{equation}
C_P = \left( \frac{\partial H}{\partial T} \right)_P
\end{equation}
donde:
\begin{itemize}
    \item $C_V$: Capacidad calorífica a volumen constante [J/K]
    \item $C_P$: Capacidad calorífica a presión constante [J/K]
    \item $H$: Entalpía del sistema [J]
    \item $T$: Temperatura [K]
\end{itemize}

El estudio del Primer Principio de la Termodinámica es esencial para comprender cómo se intercambia la energía en los sistemas físicos y cómo se puede aplicar este conocimiento en ingeniería y otras ciencias aplicadas. La capacidad para analizar y predecir los cambios de energía en sistemas termodinámicos tiene aplicaciones en una amplia variedad de campos, desde la ingeniería de procesos hasta la climatología.

\actividad{Responder el siguiente cuestionario}
\begin{enumerate}
    \item ¿Qué establece el Primer Principio de la Termodinámica?
    \item ¿Qué representa la energía interna ($U$) en un sistema termodinámico?
    \item ¿Cuál es la diferencia entre el calor ($Q$) y el trabajo ($W$) en términos de transferencia de energía?
    \item ¿Cómo se expresa matemáticamente el Primer Principio de la Termodinámica?
    \item ¿Qué es la capacidad calorífica a volumen constante ($C_V$)?
    \item ¿Qué es la capacidad calorífica a presión constante ($C_P$)?
    \item ¿Qué unidades se utilizan para medir el calor y el trabajo en termodinámica?
    \item ¿Cómo se relaciona el Primer Principio de la Termodinámica con la conservación de la energía?
    \item ¿Cuál es la importancia del Primer Principio de la Termodinámica en ingeniería?
    \item ¿Qué significa que un proceso sea adiabático en términos del Primer Principio de la Termodinámica?

\end{enumerate}

\actividad{Resolver los siguientes problemas}
\begin{enumerate}
    \item Un sistema recibe 500 J de calor y realiza 200 J de trabajo. ¿Cuál es el cambio en la energía interna del sistema?
    \item Si un gas ideal con una capacidad calorífica a volumen constante de $C_V = 20 \ \text{J/K}$ se calienta desde 300 K hasta 400 K sin realizar trabajo, ¿cuánto calor se le añade al sistema?
    \item Un sistema cerrado realiza un trabajo de 100 J y su energía interna disminuye en 50 J. ¿Cuánto calor ha intercambiado con su entorno?
    \item En un proceso isobárico, un gas ideal absorbe 1000 J de calor. Si la capacidad calorífica a presión constante del gas es $C_P = 30 \ \text{J/K}$, ¿cuál es el cambio de temperatura del gas?
    \item Un sistema tiene una energía interna inicial de 1500 J. Si se le añade 300 J de calor y se realiza un trabajo de 200 J sobre el sistema, ¿cuál es la energía interna final del sistema?
    \item Durante un proceso adiabático, un gas realiza 250 J de trabajo. ¿Cuál es el cambio en la energía interna del gas?
    \item Un gas ideal con $C_V = 25 \ \text{J/K}$ se enfría de 500 K a 300 K a volumen constante. ¿Cuál es el cambio en la energía interna del gas?
    \item Si un sistema realiza 400 J de trabajo y su energía interna aumenta en 100 J, ¿cuánto calor ha recibido?
    \item Un gas ideal se expande y realiza un trabajo de 150 J mientras recibe 200 J de calor. ¿Cuál es el cambio en la energía interna del gas?
    \item En un ciclo termodinámico completo, un sistema realiza 500 J de trabajo y recibe 500 J de calor. ¿Cuál es el cambio en la energía interna del sistema?

\end{enumerate}

\end{document}
